\QHeading{第1题　在\(\bigcirc\)和\(\square\)里分别填写相应的运算符号和数。}
\par\noindent


\PlainItem{\(25\times36=25\times4\){\color{answerred}\(\times\)}{\color{answerred}9}}
\par

\PlainItem{\(360\div8\div5=360\){\color{answerred}\(\div\)}\((8\times\){\color{answerred}5}\()\)}
\PlainItem{\(2700\div3\div9=2700\div(\){\color{answerred}3}\(\times\){\color{answerred}9}\()\)}
\PlainItem{\(150\div25\div2=\){\color{answerred}150}\(\div(\){\color{answerred}25}\(\times\){\color{answerred}2}\()\)}
\par\noindent

\QHeading{第2题　用简便方法计算下面各题。}

\CalcBlock[方法一：简便计算]{1}{\begin{aligned}[t]&\phantom{={}}\color{black}37\times5\times2\\&\color{answerred}=37\times(5\times2)\\&\color{answerred}=37\times10\\&\color{answerred}=370\end{aligned}}
\par
\MethodBlock{方法二：按运算顺序直接计算（对照）}{本法只作结果核对；本题要求简便计算，前面的简便解法仍保留。}{\begin{aligned}[t]&\phantom{={}}37\times 5\times 2\\
&=185\times 2\\
&=370\end{aligned}}


\CalcBlock[方法一：简便计算]{2}{\begin{aligned}[t]&\phantom{={}}\color{black}28\times2\times50\\&\color{answerred}=28\times(2\times50)\\&\color{answerred}=28\times100\\&\color{answerred}=2800\end{aligned}}
\MethodBlock{方法二：结合因数，先凑整}{用乘法交换律、结合律，先算容易凑整的两个因数。}{\begin{aligned}[t]&\phantom{={}}28\times 2\times 50\\
&=(28\times 50)\times 2\\
&=1400\times 2\\
&=2800\end{aligned}}

\MethodBlock{方法三：先处理括号或乘除，再把特殊因数改成整百或整千}{先把25、125等改写成整百或整千除以一个数。}{\begin{aligned}[t]&\phantom{={}}28\times 2\times 50\\
&=56\times 50\\
&=(100\div 2)\times 56\\
&=100\times (56\div 2)\\
&=100\times 28\\
&=2800\end{aligned}}

\CalcBlock[方法一：简便计算]{3}{\begin{aligned}[t]&\phantom{={}}\color{black}125\times7\times8\\&\color{answerred}=(125\times8)\times7\\&\color{answerred}=1000\times7\\&\color{answerred}=7000\end{aligned}}
\MethodBlock{方法二：先处理括号或乘除，再凑整后补偿}{把接近整十、整百的因数凑整，再补回或减去多算的部分。}{\begin{aligned}[t]&\phantom{={}}125\times 7\times 8\\
&=875\times 8\\
&=875\times (10-2)\\
&=875\times 10-875\times 2\\
&=8750-1750\\
&=7000\end{aligned}}

\CalcBlock[方法一：简便计算]{4}{\begin{aligned}[t]&\phantom{={}}\color{black}25\times14\times40\\&\color{answerred}=(25\times40)\times14\\&\color{answerred}=1000\times14\\&\color{answerred}=14000\end{aligned}}
\MethodBlock{方法二：结合因数，先凑整}{用乘法交换律、结合律，先算容易凑整的两个因数。}{\begin{aligned}[t]&\phantom{={}}25\times 14\times 40\\
&=(14\times 40)\times 25\\
&=560\times 25\\
&=14000\end{aligned}}

\MethodBlock{方法三：按运算顺序直接计算（对照）}{本法只作结果核对；本题要求简便计算，前面的简便解法仍保留。}{\begin{aligned}[t]&\phantom{={}}25\times 14\times 40\\
&=350\times 40\\
&=14000\end{aligned}}

\CalcBlock[方法一：简便计算]{5}{\begin{aligned}[t]&\phantom{={}}\color{black}125\times47\times8\\&\color{answerred}=(125\times8)\times47\\&\color{answerred}=1000\times47\\&\color{answerred}=47000\end{aligned}}
\MethodBlock{方法二：先处理括号或乘除，再凑整后补偿}{把接近整十、整百的因数凑整，再补回或减去多算的部分。}{\begin{aligned}[t]&\phantom{={}}125\times 47\times 8\\
&=5875\times 8\\
&=5875\times (10-2)\\
&=5875\times 10-5875\times 2\\
&=58750-11750\\
&=47000\end{aligned}}

\CalcBlock[方法一：简便计算]{6}{\begin{aligned}[t]&\phantom{={}}\color{black}5\times69\times2\\&\color{answerred}=(5\times2)\times69\\&\color{answerred}=10\times69\\&\color{answerred}=690\end{aligned}}
\MethodBlock{方法二：按运算顺序直接计算（对照）}{本法只作结果核对；本题要求简便计算，前面的简便解法仍保留。}{\begin{aligned}[t]&\phantom{={}}5\times 69\times 2\\
&=345\times 2\\
&=690\end{aligned}}

\par\noindent

\QHeading{第3题　解决问题。}

\Problem{①商店运来4箱牛奶，每箱有25盒，每盒售价2元。全部售出后，一共可以卖多少元？}
\par


\MethodTitle{方法一}
\Answer{\Formula{\begin{aligned}[t]&\phantom{={}}2\times25\times4\\&=2\times(25\times4)\\&=2\times100\\&=200\text{（元）}\end{aligned}}\par 答：全部售出后，一共可以卖200元。}

\MethodTitle{方法二：先求一箱能卖的钱}
\Answer{每盒2元，每箱25盒，先求每箱的总价，再乘箱数。}
\Answer{\Formula{\begin{aligned}[t]&\phantom{={}}2\times 25\times 4\\
&=50\times 4\\
&=200\text{（元）}\end{aligned}}}
\Answer{答：全部售出一共可以卖200元。}
\Problem{②学校买来8捆跳绳，每捆125根，每根3元。买这些跳绳一共花了多少钱？}

\MethodTitle{方法一}
\Answer{\Formula{\begin{aligned}[t]&\phantom{={}}3\times125\times8\\&=3\times(125\times8)\\&=3\times1000\\&=3000\text{（元）}\end{aligned}}\par 答：买这些跳绳一共花了3000元。}

\MethodTitle{方法二：先求每捆跳绳的价钱}
\Answer{每捆125根，每根3元，先求每捆总价，再算8捆。}
\Answer{\Formula{\begin{aligned}[t]&\phantom{={}}3\times 125\times 8\\
&=375\times 8\\
&=3000\text{（元）}\end{aligned}}}
\Answer{答：一共花了3000元。}
